\documentclass[12pt]{article}

\usepackage[margin=1in]{geometry}
\usepackage{setspace}
\usepackage{graphicx}
\usepackage{hyperref}
\usepackage{enumitem}

\title{Project Title Here}
\author{Group Member 1 \and Group Member 2 \and Group Member 3}
\date{\today}

\begin{document}
\maketitle
\onehalfspacing

\section{Introduction}
\textbf{Project Question/Goal:} TODO.

TODO: Describe the problem domain, motivation, and required background context.
Update your Phase 1 proposal content based on TA feedback and any direction changes.
Add inline citations for facts, datasets, and prior work.

\section{Datasets}
\textit{If your project does not use datasets, state that clearly and explain your data source/generation approach.}

\subsection{Dataset 1}
\begin{itemize}[leftmargin=*]
    \item Name: TODO
    \item Source: TODO (include URL and access date)
    \item Format: TODO (e.g., CSV/JSON)
    \item Fields used by program: TODO
    \item Notes on preprocessing/subsetting: TODO
\end{itemize}

\subsection{Dataset 2}
\begin{itemize}[leftmargin=*]
    \item Name: TODO
    \item Source: TODO
    \item Format: TODO
    \item Fields used by program: TODO
    \item Notes on preprocessing/subsetting: TODO
\end{itemize}

\section{Computational Overview}
TODO: Describe the actual program implementation at a high level.

\subsection{Data Representation (Trees/Graphs)}
TODO: Explain how your data structures model the domain and where trees/graphs are central.

\subsection{Major Computations and Algorithms}
TODO: Summarize key computations (loading/building structures, filtering/aggregation, algorithms/models).
Reference concrete files, classes, and core functions.

\subsection{Output, Visualization, and Interactivity}
TODO: Explain what output users see and how they interact with it (if applicable).

\subsection{External Python Libraries}
TODO: List new libraries and explain the relevant functions/data types used.

\section{Instructions: Dataset Access and Running the Program}
\subsection{Python and Library Setup}
\begin{enumerate}[leftmargin=*]
    \item Install Python 3.13.5.
    \item Install dependencies from \texttt{requirements.txt}.
    \item TODO: Add any extra installation notes if needed.
\end{enumerate}

\subsection{Dataset Download Instructions}
Choose one method and provide exact steps:
\begin{itemize}[leftmargin=*]
    \item MarkUs upload (zip under 10 MB), or
    \item \url{https://send.utoronto.ca} shared with \texttt{csc111-2026-01@cs.toronto.edu}, or
    \item U of T OneDrive public/UTORid-accessible link.
\end{itemize}
TODO: Add exact link/location and extraction instructions.

\subsection{How to Run}
\begin{enumerate}[leftmargin=*]
    \item Place all Python files in one folder (no subfolders for modules).
    \item Ensure datasets are extracted to: TODO path.
    \item Run \texttt{main.py}.
\end{enumerate}

Expected output:
\begin{itemize}[leftmargin=*]
    \item TODO: Describe what appears when \texttt{main.py} runs.
    \item TODO: Describe interaction steps and what the TA should try.
    \item TODO: Add screenshots with \texttt{\textbackslash includegraphics} if helpful.
\end{itemize}

\section{Changes from Proposal to Final Submission}
TODO: Explain key scope/direction/design changes and why they were made.
Mention TA/instructor feedback and any technical constraints that influenced decisions.

\section{Discussion}
TODO: Analyze whether results answer your question/goal.
Discuss limitations (data quality, algorithms, tools, runtime, assumptions) and next steps.
Target approximately 500--800 words.

\section{References}
% You can replace this with a BibTeX workflow if preferred.
\begin{thebibliography}{9}
\bibitem{ref1}
TODO: Topic research source.

\bibitem{ref2}
TODO: Dataset source.

\bibitem{ref3}
TODO: Library documentation/tutorial source.
\end{thebibliography}

\end{document}
